% ---------------------------------------------------------------------
% 2.14 Estado del arte
% ---------------------------------------------------------------------
\section[ESTADO DEL ARTE]{Estado del arte}

\subsection{Sistemas P300 Speller y competencias BCI}

A partir del deletreador de Farwell y Donchin \cite{farwell1988}, la investigación sobre el P300
Speller se centró durante varios años en mejorar la clasificación de las épocas con métodos
lineales. Las competencias internacionales de BCI jugaron un papel central en este proceso, pues
pusieron a disposición de la comunidad bases de datos comunes; en la tercera de ellas, el conjunto
de SVM lineales propuesto por Rakotomamonjy y Guigue obtuvo el primer lugar al identificar
correctamente el 96.5\% de los caracteres de prueba con 15 repeticiones \cite{rakotomamonjy2008}.
En paralelo, Krusienski et al. mostraron que el SWLDA y el discriminante de Fisher eran los
clasificadores más robustos entre los evaluados \cite{krusienski2006}, y posteriormente que una
adecuada selección de ocho electrodos permitía mejorar el desempeño respecto al uso exclusivo de la
línea media \cite{krusienski2008}.

\subsection{Sistemas de comunicación con usuarios con discapacidad}

La validación con usuarios reales ha sido otra línea importante de investigación. Sellers y
Donchin realizaron las primeras pruebas de un P300 Speller con pacientes con ELA
\cite{sellers2006}, y Hoffmann et al. propusieron un sistema de seis opciones evaluado con personas
con distintas discapacidades motoras, cuyos registros constituyen hoy una de las bases de datos
públicas más utilizadas \cite{hoffmann2008}. En el ámbito latinoamericano, el estudio ``Interfaz
cerebro-computador basada en P300 para la comunicación alternativa'' \cite{garciacossio2011}
implementó un sistema funcional de deletreo utilizando una matriz de estimulación visual con dos
adolescentes en situación de discapacidad motora, reportando porcentajes de acierto promedio de 95\%
y 85\% para cada usuario en la prueba de libre deletreo. Por su parte, el trabajo de Brunner et
al. \cite{brunner2011} demostró la operación de un sistema P300 Speller utilizando señales de
electrocorticografía, alcanzando tasas sostenidas superiores a 17 caracteres por minuto; no
obstante, el uso de técnicas invasivas limita su aplicación en contextos clínicos generales, donde
el EEG continúa siendo la alternativa más accesible.

\subsection{Aprendizaje profundo para la detección de la P300}

Con la introducción de las redes convolucionales, Cecotti y Gräser mostraron que era posible
aprender directamente los filtros espaciales y temporales a partir de la señal, obteniendo en la
competencia BCI III resultados comparables a los del ganador \cite{cecotti2011}. A partir de
entonces se propusieron arquitecturas como BN3 \cite{liu2018}, OCLNN \cite{shan2018} y EEGNet
\cite{lawhern2018}, y la comparación de Alvarado-González et al. mostró que redes muy compactas
como SepConv1D, con apenas cuatro filtros en su única capa convolucional, pueden alcanzar un AUC
similar a la de redes mucho más grandes \cite{alvarado2021}.

En el ámbito nacional, la tesis de maestría titulada ``Sistema clasificador de señales EEG para
el control de una neuroprótesis por medio de un dispositivo móvil'' \cite{contreras2021},
desarrollada en la Escuela Superior de Cómputo del Instituto Politécnico Nacional, implementó un
sistema basado en la arquitectura SepConv1D para la detección del componente P300 con solo cuatro
electrodos secos, reportando valores promedio de AUC cercanos a 0.82 en evaluaciones intrasujeto y
de 0.70 en escenarios intersujeto sin calibración individual al considerar filas y columnas. Este
trabajo demostró la viabilidad técnica de integrar procesamiento EEG, extracción automática de
características y modelos de aprendizaje profundo en un sistema funcional; sin embargo, su
aplicación estuvo orientada al control de una neuroprótesis y no específicamente a la
implementación de un sistema de comunicación mediante deletreo. En la misma escuela, Santos,
Romero y Fernández desarrollaron una aplicación para manipular un brazo robótico virtual mediante
potenciales P300 con matrices de $5\times5$ de iconos y de letras, clasificadas con el
discriminante lineal de Fisher \cite{santos2017}, y Anguiano y Ramírez utilizaron señales de EEG
y redes neuronales artificiales para detectar estrés en estudiantes \cite{anguiano2019}.

\subsection{Modelos de lenguaje en el P300 Speller}

La incorporación de información lingüística es la línea más reciente para acelerar la escritura.
Ryan et al. \cite{ryan2011} y Speier et al. \cite{speier2012} mostraron que la predicción de
palabras y los modelos de lenguaje reducen el tiempo de escritura o mejoran la tasa de
transferencia, y más recientemente la investigación ``ChatBCI'' \cite{hong2026} incorporó modelos
de lenguaje de gran escala para mejorar las selecciones dentro de un sistema P300 Speller, logrando
una reducción del tiempo promedio de escritura de 62.14\%.

\subsection{Resumen comparativo}

De los trabajos revisados se identifican dos limitaciones persistentes en los sistemas P300
basados en EEG: la necesidad de múltiples canales para mantener niveles altos de precisión y la
dependencia de múltiples repeticiones de estímulos para asegurar una detección confiable. La
Tabla~\ref{tab:estado-arte} resume los enfoques descritos, contrastando el número de canales
empleados, el clasificador, el uso de un modelo de lenguaje y el desempeño reportado por cada
trabajo frente a la propuesta del presente Trabajo Terminal.

\begin{table}[htbp]
\centering
\caption{Resumen comparativo del estado del arte}
\label{tab:estado-arte}
\footnotesize
\begin{tabularx}{\textwidth}{L{3.1cm}L{1.8cm}L{2.4cm}L{2.1cm}Y}
\toprule
\textbf{Trabajo} & \textbf{Señal / canales} & \textbf{Clasificador} & \textbf{Modelo de lenguaje} & \textbf{Desempeño} \\
\midrule
Farwell y Donchin, 1988 \cite{farwell1988} & EEG / 1 (Pz) & SWDA, pico, área y covarianza & No & Exactitud $\approx$ 95\%; 2.3 caracteres/min \\
Sellers y Donchin, 2006 \cite{sellers2006} & EEG & -- & No & Primeras pruebas con pacientes con ELA \\
Rakotomamonjy y Guigue, 2008 \cite{rakotomamonjy2008} & EEG / 64 & Conjunto de SVM & No & 96.5\% de caracteres con 15 repeticiones \\
Krusienski et al., 2008 \cite{krusienski2008} & EEG / 8 & SWLDA & No & Conjunto de 8 canales superior a la línea media \\
Cecotti y Gräser, 2011 \cite{cecotti2011} & EEG / 64 & CNN & No & $\approx$ 95\% de caracteres con 15 repeticiones \\
García Cossio et al., 2011 \cite{garciacossio2011} & EEG / 6 & -- & No & Acierto de 95\% y 85\% en libre deletreo \\
Brunner et al., 2011 \cite{brunner2011} & ECoG / 96 subdurales & SWLDA & No & 17 caracteres/min \\
Ryan et al., 2011 \cite{ryan2011} & EEG & -- & Diccionario & Menor tiempo por frase \\
Speier et al., 2012 \cite{speier2012} & EEG & Clasificación dinámica & Probabilístico & Mayor exactitud y tasa de transferencia \\
Santos et al., 2017 (ESCOM) \cite{santos2017} & EEG & LDA de Fisher & No & Precisión de 90.01\% (20 sujetos) \\
Alvarado-González et al., 2021 \cite{alvarado2021} & EEG / varios & SepConv1D y otras 11 redes & No & AUC intra $\approx$ 0.84; inter $\approx$ 0.78 \\
Contreras, 2021 (IPN) \cite{contreras2021} & EEG / 4 & SepConv1D & No & AUC $\approx$ 0.82 (intrasujeto), $\approx$ 0.70 (intersujeto) \\
ChatBCI, Hong et al., 2026 \cite{hong2026} & EEG / 32 & -- & LLM & Reducción de tiempo promedio de 62.14\% \\
\midrule
\textbf{Propuesta TT} & EEG / 4 & Aprendizaje profundo (CNN) & Predicción de palabras & -- \\
\bottomrule
\end{tabularx}
\par\smallskip{\footnotesize Elaboración propia.}
\end{table}
